\documentclass[11pt,a4paper]{article}
\usepackage[utf8]{inputenc}
\usepackage[T1]{fontenc}
\usepackage[margin=1in]{geometry}
\usepackage[hidelinks]{hyperref}
\usepackage{enumitem}

\begin{document}

\begin{center}
    {\Huge \textbf{Thiago Gomes Stein}} \\
    \vspace{2mm}
    Londrina, PR | +55 (43) 99991-5412 | \href{mailto:thiagomes.stein@gmail.com}{thiagomes.stein@gmail.com} \\
    \vspace{1mm}
    \href{https://www.linkedin.com/in/thiago-gomes-stein-31649a35b/}{LinkedIn} | \href{https://github.com/Thiago-Stein}{GitHub} | \href{https://lattes.cnpq.br/0305313799581622}{Currículo Lattes}
\end{center}

\vspace{4mm}

\section*{Resumo Profissional}
Estudante de Engenharia de Software focado em desenvolvimento Backend. Possuo experiência prática com Java, ecossistema JavaScript (Node.js) e bancos de dados relacionais. Usuário avançado de sistemas Linux (Arch Linux e Mint), construindo aplicações eficientes. Me destaco pelo autodidatismo, curiosidade em aprender, disciplina e forte capacidade de comunicação, habilidades desenvolvidas através da participação ativa em grupos de estudo intensivos e na atuação contínua como monitor acadêmico de tecnologia.

\section*{Formação Acadêmica}
\textbf{Engenharia de Software} -- Bacharelado \\
Centro Universitário Filadélfia (UniFil) -- Londrina, PR \hfill Fev. 2025 -- Previsão: Dez. 2028

\textbf{Ensino Médio} \\
Colégio Estadual Polivalente de Londrina -- Londrina, PR \hfill 2021 -- 2023

\section*{Experiência Acadêmica e Extracurricular}

\textbf{Membro do Núcleo de Prática em Informática (NPI)} | \textit{UniFil} \hfill Jan. 2026 -- Presente
\begin{itemize}[leftmargin=*]
    \item Participação diária em grupo de estudos autodidata (20 horas semanais) focado no aprofundamento prático em tecnologia e programação.
    \item Desenvolvimento de forte disciplina, autogestão e aprendizado contínuo, conciliando rigoroso controle de metas e presença com as atividades de monitoria acadêmica.
\end{itemize}

\vspace{2mm}

\textbf{Monitor de Tecnologia} | \textit{Projeto Londrinense Tech (UniFil)} \hfill Fev. 2025 -- Presente
\begin{itemize}[leftmargin=*]
    \item \textbf{Atuação Atual (2026 - Presente):} Mentoria para turmas do Ensino Fundamental com foco em eletrônica básica e prototipagem, prestando suporte prático na utilização de Tinkercad e Arduino.
    \item \textbf{Atuação Anterior (2025):} Suporte técnico e pedagógico para turmas do Ensino Fundamental e Médio nas disciplinas de Java, Introdução a Banco de Dados e Lógica de Programação (Portugol).
    \item Auxílio direto aos professores titulares no andamento das aulas e esclarecimento de dúvidas individuais.
\end{itemize}

\section*{Projetos Principais}

\textbf{API Filmes} | \href{https://github.com/Thiago-Stein/trabalhoFinalAPI}{Ver no GitHub} \hfill \textit{Node.js, JavaScript, SQLite 3, Postman}
\begin{itemize}[leftmargin=*]
    \item Desenvolvimento de uma API RESTful completa para gerenciamento de listas de filmes.
    \item Estruturação do banco de dados relacional e realização de testes rigorosos de rotas utilizando o Postman.
\end{itemize}

\textbf{Projeto Perguntalys} | \href{https://github.com/Thiago-Stein/jogo-perguntalys}{Ver no GitHub} | \href{https://thiago-stein.github.io/jogo-perguntalys/}{Ver jogo publicado} \hfill \textit{HTML5, CSS, JavaScript}
\begin{itemize}[leftmargin=*]
    \item Jogo web educacional focado no desenvolvimento de raciocínio lógico para o público infantojuvenil.
    \item Desenvolvimento colaborativo em equipe ágil, atuando na programação front-end e gerenciamento de estado via persistência local (LocalStorage).
\end{itemize}

\section*{Habilidades Técnicas}
\begin{itemize}[leftmargin=*]
    \item \textbf{Linguagens:} Java, JavaScript, SQL, HTML5, CSS3, Portugol, C/C++ (Arduino).
    \item \textbf{Backend \& Ferramentas:} Node.js, Postman, Git, GitHub.
    \item \textbf{Bancos de Dados:} SQLite, Modelagem de Dados (DDL, Formas Normais, Cardinalidade).
    \item \textbf{Infraestrutura \& Outros:} Linux (Arch Linux e Mint), Tinkercad.
\end{itemize}

\section*{Idiomas}
\begin{itemize}[leftmargin=*]
    \item \textbf{Inglês:} Intermediário. \hfill \textbf{Português:} Nativo.
\end{itemize}

\end{document}
